The continuous-time market model enables the incorporation of portfolio reallocation algorithms within a stochastic dynamic programming framework.

\subsection{Asset Price Modelling}
\noindent The prices at time $t \ge 0$ of $d+1$ assets numbered $0, 1, \dots, d$ is denoted by the random vector
\begin{equation} \label{eq:SpotPriceVector}
    \xbar{S}_t = \left(S_t^{(0)}, S_t^{(1)}, \dots, S_t^{(d)}\right)
\end{equation}
which forms a stochastic process $(\xbar{S}_t)_{t \in \R_+}$. As in discrete time, the asset $n^o$ 0 is a riskless asset (of savings account type) yielding an interest rate $r$, \textit{i.e.} we have
\begin{equation} \label{eq:RisklessAssetContinuous}
    S_t^{(0)} = S_0^{(0)} \e^{rt}, \quad t \ge 0.
\end{equation}

\begin{defn}{Discounting}{DiscountingDef}
    Let
    \begin{equation} \label{eq:DiscountedVectorDef}
        \xbar{X}_t \coloneqq \left(\widetilde{S}_t^{(0)}, \widetilde{S}_t^{(1)}, \dots, \widetilde{S}_t^{(d)}\right), \quad t \in \R,
    \end{equation}
    denote the vector of discounted asset prices, defined as:
    \begin{equation} \label{eq:DiscountedComponentDef}
        \widetilde{S}_t^{(k)} = \e^{-rt}S_t^{(k)}, \quad t \ge 0, \quad k = 0, 1, \dots, d.
    \end{equation}
\end{defn}

\noindent We can also write
\begin{equation} \label{eq:DiscountedVectorCompact}
    \xbar{X}_t \coloneqq \e^{-rt}\xbar{S}_t, \quad t \ge 0.
\end{equation}

\begin{defn}{Portfolio Strategy}{PortfolioStrategyDef}
    A portfolio strategy is a stochastic process $(\xbar{\xi}_t)_{t \in \R_+} \subset \R^{d+1}$, where $\xi_t^{(k)}$ denotes the (possibly fractional) quantity of asset $n^o$ $k$ held at time $t \ge 0$.
\end{defn}

\noindent The value at time $t \ge 0$ of the portfolio strategy $(\xbar{\xi}_t)_{t \in \R_+} \subset \R^{d+1}$ is defined by
\begin{equation} \label{eq:PortfolioValueDef}
    V_t \coloneqq \xbar{\xi}_t \cdot \xbar{S}_t = \sum_{k=0}^d \xi_t^{(k)} S_t^{(k)}, \quad t \ge 0.
\end{equation}

\noindent The discounted value of the portfolio is defined by
\begin{align}
    \widetilde{V}_t &\coloneqq \e^{-rt} V_t \label{eq:DiscountedValueDef} \\
    &= \e^{-rt} \xbar{\xi}_t \cdot \xbar{S}_t \nonumber \\
    &= \e^{-rt} \sum_{k=0}^d \xi_t^{(k)} S_t^{(k)} \nonumber \\
    &= \sum_{k=0}^d \xi_t^{(k)} \widetilde{S}_t^{(k)} \nonumber \\
    &= \xbar{\xi}_t \cdot \xbar{X}_t, \quad t \ge 0. \label{eq:DiscountedValueCompact}
\end{align}

\noindent The effect of discounting from time $t$ to time $0$ is to divide prices by $\e^{rt}$, making all prices comparable at time $0$.

\subsection{Arbitrage and Risk-Neutral Measures}
\noindent In continuous-time, the definition of arbitrage follows the lines of its analogs in the one-step and discrete-time models. In what follows we will only consider \textit{admissible} portfolio strategies whose total value $V_t$ remains nonnegative for all times $t \in [0, T]$.

\begin{defn}{Admissible Portfolio Strategy}{AdmissiblePortfolioStrategyDef}
    A portfolio strategy $(\xi_t^{(k)})_{t \in [0, T], k=0, 1, \dots, d}$ with value
    \begin{equation} \label{eq:AdmissiblePortfolioValue}
        V_t = \xbar{\xi}_t \cdot \xbar{S}_t = \sum_{k=0}^d \xi_t^{(k)} S_t^{(k)}, \quad t \in [0, T],
    \end{equation}
    constitutes an arbitrage opportunity if all three following conditions are satisfied:
    \begin{enumerate}[label=\roman*), leftmargin=*, parsep=6pt]
        \item $V_0 \le 0$ at time $t=0$, \hfill [Start from a zero-cost portfolio, or with a debt.]
        \item $V_T \ge 0$ at time $t=T$, \hfill [Finish with a nonnegative amount.]
        \item $\Pm(V_T > 0) > 0$ at time $t=T$. \hfill [Profit is made with nonzero probability.]
    \end{enumerate}
\end{defn}

\noindent Roughly speaking, $(ii)$ means that the investor wants no loss, $(iii)$ means that he wishes to sometimes make a strictly positive gain, and $(i)$ means that he starts with zero capital or even with a debt.

Next, we turn to the definition of risk-neutral probability measures (or martingale measures) in continuous time, which states that under a risk-neutral probability measure $\Pm^*$, the return of the risky asset over the time interval $[u, t]$ equals the return of the riskless asset given by
\begin{equation} \label{eq:RisklessReturnContinuous}
    S_t^{(0)} = \e^{(t-u)r} S_u^{(0)}, \quad 0 \le u \le t.
\end{equation}
Recall that the filtration $(\cF_t)_{t \in \R_+}$ is generated by Brownian motion $(B_t)_{t \in \R_+}$, \textit{i.e.}
\begin{equation} \label{eq:BrownianFiltration}
    \cF_t = \sigma(B_u : 0 \le u \le t), \quad t \ge 0.
\end{equation}

\subsection{Self-financing Portfolio Strategies}
Let $\xi_t^{(i)}$ denote the (possibly fractional) quantity invested at time $t$ over the time interval $[t, t + \di t)$, in the asset $S_t^{(k)}$, $k = 0, 1, \dots, d$, and let
\begin{equation}
    \xbar{\xi}_t = \left(\xi_t^{(k)}\right)_{k=0,1,\dots,d}, \quad \xbar{S}_t = \left(S_t^{(k)}\right)_{k=0,1,\dots,d}, \quad t \ge 0,
\end{equation}
denote the associated portfolio value and asset price processes. The portfolio value $V_t$ at time $t \ge 0$ is given by
\begin{equation}
    V_t = \xbar{\xi}_t \cdot \xbar{S}_t = \sum_{k=0}^d \xi_t^{(k)} S_t^{(k)}, \quad t \ge 0.
    \label{eq:PortfolioValueContinuous}
\end{equation}

\begin{defn}{Self-Financing Portfolio Update}{SelfFinancingPortfolioUpdate}
    The portfolio strategy $(\xbar{\xi}_t)_{t \in \R_+}$ is self-financing if the portfolio value remains constant after updating the portfolio from $\xbar{\xi}_t$ to $\xbar{\xi}_{t+\di t}$, i.e.
    \begin{equation}
        \xbar{\xi}_t \cdot \xbar{S}_{t+\di t} = \sum_{k=0}^d \xi_t^{(k)} S_{t+\di t}^{(k)} = \sum_{k=0}^d \xi_{t+\di t}^{(k)} S_{t+\di t}^{(k)} = \xbar{\xi}_{t+\di t} \cdot \xbar{S}_{t+\di t},
        \label{eq:SelfFinancingConditionContinuous}
    \end{equation}
    which is the continuous-time analog of the self-financing condition already encountered in the discrete setting of \cref{sec:AssetsPortfoliosArbitrage}, see \ref{defn:PortfolioDef}. A major difference with the discrete-time case of \ref{defn:PortfolioDef}, however, is that the continuous-time differentials $\di S_t$ and $\di \xi_t$ do not make pathwise sense as continuous-time stochastic integrals are defined by $L^2$ limits, cf. \ref{prop:CauchyCriteriaLp}, or by convergence in probability.
\end{defn}